% =====================================================================
%  CHAPTER 6: 儿童心理问题（对应大纲第六章）
%  内容来源：往届笔记第十二章，参考PPT第十章
% =====================================================================
\chapter{儿童心理问题}
\setcounter{chapter}{6}
\chapterintro{
\textbf{掌握：}心理卫生的基本概念、心理健康标准和工作目标，童年期和青春期的心理特点和常见的心理卫生问题。
}

% ===== 掌握内容 =====
\section{心理卫生概述}

\subsection{基本概念}

\begin{defbox}
\textbf{心理卫生（mental health）/ 精神卫生：}研究如何维护和促进人心理健康的科学。

\textbf{心理障碍（mental disorder）：}儿童少年在心理健康方面存在的偏倚统称心理卫生问题；若其严重程度、持续时间超过相应年龄的允许范围为心理障碍。
\end{defbox}

\subsection{儿童少年心理健康标准}

\begin{keybox}
\textbf{儿童少年心理健康标准：}

\begin{enumerate}[leftmargin=*]
    \item \imp{智力发展正常}
    \item \imp{情绪稳定而反应适度}
    \item \imp{心理行为特点符合年龄}
    \item \imp{人际关系的心理适应}——能与人和睦相处，悦纳自己，认同他人
    \item \imp{个性的稳定和健全}——表现出健康的精神风貌，客观的自我意识，行为符合道德规范，能适度耐受各种压力和应激
\end{enumerate}
\end{keybox}

\subsection{儿童少年心理卫生特点与工作目标}

\begin{keybox}
\textbf{儿童少年心理卫生特点：}

\begin{enumerate}[leftmargin=*]
    \item 不同年龄阶段的人群有各自的生理及心理特点，因此心理卫生的内容也不相同。
    \item 儿童少年心理卫生不仅与个体认知、情绪发展相关，更受家庭、学校及社会的影响。
    \item 危险因素较多、累积程度较高时，可能引发心理卫生问题，甚至造成心理障碍。
    \item 保护因素较多时，更有可能维持心理健康的动态平衡。
\end{enumerate}

\textbf{工作目标：}

维护和促进儿童少年心理健康，预防心理障碍的发生，早期发现和干预心理卫生问题，促进儿童少年全面发展。
\end{keybox}

\subsection{儿童少年心理障碍病因}

\begin{tipbox}
\textbf{病因特点：}

\begin{enumerate}[leftmargin=*]
    \item 儿童少年心理发展的特点和异常行为有多样性，优势和不足常共存，很多行为问题或障碍不是由简单清晰的因果关系导致。
    \item 儿童少年心理行为障碍的病因复杂多样，同样的心理障碍可能表现形式不同。
    \item 导致特定障碍的途径是多样交互的，而非线性静态的。
\end{enumerate}
\end{tipbox}

\section{童年期和青春期的心理特点}\difficulty

\begin{keybox}
\textbf{一、童年期心理特点}

\begin{enumerate}[leftmargin=*]
    \item \textbf{认知发展：}思维从具体形象思维逐步过渡到抽象逻辑思维；有意注意进一步发展；记忆发生质的变化。
    \item \textbf{情绪发展：}情感内容逐渐丰富；情感的深刻性和稳定性逐渐增加；高级情感逐渐发展。
    \item \textbf{社会性发展：}学校集体生活成为重要的社会化途径；同伴关系日益重要；自我意识逐步发展。
\end{enumerate}

\textbf{二、青春期心理特点}

\begin{enumerate}[leftmargin=*]
    \item \textbf{认知发展：}抽象逻辑思维处于优势地位；记忆整体水平处于人生最佳时期；有意注意稳定约40min。
    \item \textbf{情绪发展：}情绪情感迅速而热烈，有两极性；内隐性和表现性共同存在；高级情感已有相当发展。
    \item \textbf{自我意识发展：}自我意识迅猛发展；出现"心理断乳"现象；追求独立自主。
    \item \textbf{性心理发展：}经历异性疏远期→异性爱慕期→两性初恋期；存在性生理成熟与性心理幼稚的矛盾。
    \item \textbf{社会性发展：}同伴影响力增强；价值观和人生观逐步形成；容易出现逆反心理。
\end{enumerate}

△青春期是性心理障碍开始表现的关键期。
\end{keybox}

\section{儿童少年常见心理卫生问题}

\subsection{注意缺陷多动障碍（ADHD）}

\begin{keybox}
\textbf{注意缺陷多动障碍（attention-deficit/hyperactivity disorder, ADHD）/ 多动症：}注意缺陷在行为上表现为游离于任务、缺乏持续性、难以维持注意力以及杂乱无章，且并非由于违拗或缺乏理解力所致。多动是指在不恰当的时候过多的躯体运动（例如孩子到处乱跑），或过于坐立不安、手脚动个不停或讲话过多（DSM-V）。起病于学龄期。

\textbf{表现特征：}

\begin{enumerate}[leftmargin=*]
    \item \textbf{过度活动：}自幼易兴奋、活动量大、多哭闹、睡眠差、喂食困难；入学课堂纪律差、无法静心作业，做事唐突、冒失。
    \item \textbf{注意力不集中：}上课时注意易被无关刺激吸引分散，无心听讲，东张西望，影响学习。
    \item \textbf{冲动：}易兴奋，做事不顾及后果；常有攻击行为；不遵守规则，缺乏忍耐，不愿等待；难以理解他人内心活动、表情，对同学的玩笑有过激反应。
    \item \textbf{学习困难、学习成绩不良：}有语言理解或表达问题，短时记忆困难等。
\end{enumerate}

☆有不良情绪（原发性情绪障碍）
\end{keybox}

\subsection{特殊性学习障碍（SLD）}

\begin{keybox}
\textbf{特殊性学习障碍（specific learning disorder, SLD）：}学龄儿童在阅读、书写、拼字、表达、计算等认知学习过程中存在一种或一种以上特殊性障碍，基本特征是对学习关键学业技能有持续性困难，分轻中重3类。

\textbf{病因学：}SLD发病原因基本与ADHD相同，但儿童母语特性也有一定关联。

\textbf{表现特征：}

\begin{enumerate}[leftmargin=*]
    \item \textbf{阅读困难：}文字阅读正确率低，朗读时易漏字或添字，读同音异义字困难或混用，阅读时多用手指指字；阅读缺乏节奏，不流畅；阅读理解能力弱，文章理解困难等。
    \item \textbf{书写表达障碍：}拼写准确性低，语法和标点符号的使用准确率低，说话缺少关系词，因果顺序表达欠佳等。书面表达清晰度或组织性差。
    \item \textbf{计算困难：}顺序或空间认知障碍，缺乏数字感，常将数字顺序颠倒，算术事实记忆不良，计算准确率或流畅性低，判断方位、距离、图形困难。
\end{enumerate}

☆无不良情绪（原发性情绪障碍）
\end{keybox}

\begin{exambox}
\textbf{△ ADHD与SLD鉴别要点：}ADHD有不良情绪（原发性情绪障碍），SLD无不良情绪。
\end{exambox}

\subsection{孤独症谱系障碍（ASD）}

\begin{keybox}
\textbf{孤独症谱系障碍（autistic-spectrum disorder, ASD）/ 自闭症：}是一种以社交沟通障碍和重复刻板行为为主要特征的发育障碍性疾病。多种场合下，社交交流和社交互动方面存在持续性缺陷。

\textbf{表现特征：}

\begin{enumerate}[leftmargin=*]
    \item \textbf{语言障碍或落后：}为主要就诊原因，具备言语能力但缺乏沟通性质。
    \item \textbf{社交障碍：}核心症状，缺乏社交技巧，生命早期即表现出目光对视交流的缺乏，不能参与合作性游戏，甚至难以与母亲建立依恋关系。
    \item \textbf{兴趣狭隘和刻板行为：}对某些物件或活动表现异常的兴趣，伴重复、刻板动作，有时表现强迫行为或自伤行为。
    \item \textbf{认知能力障碍：}大部分智力落后，部分正常或超常，岛状能力，部分有特殊才能，但在人际交往、情绪识别和调控等方面的困难。
    \item \textbf{感知觉障碍：}对疼痛不敏感，却对某些声音或触摸表现出强烈的偏好或极端的反感等。
\end{enumerate}
\end{keybox}

% ----- 复习题 -----
\reviewcontent{
\section{复习题}

\begin{enumerate}[leftmargin=*, itemsep=8pt]
    \item \textbf{【名词解释】}心理卫生；心理障碍；注意缺陷多动障碍（ADHD）；特殊性学习障碍（SLD）；孤独症谱系障碍（ASD）
    \item \textbf{【简答题】}简述儿童少年心理健康的五条标准。
    \item \textbf{【简答题】}简述童年期和青春期的心理特点。
    \item \textbf{【简答题】}简述ADHD的主要表现特征。
    \item \textbf{【简答题】}比较ADHD与SLD的主要区别。
    \item \textbf{【简答题】}简述ASD的核心症状和表现特征。
\end{enumerate}

\subsection*{参考答案要点}

\begin{enumerate}[leftmargin=*, itemsep=5pt]
    \item \textbf{心理卫生}：研究如何维护和促进人心理健康的科学。\textbf{心理障碍}：儿童少年在心理健康方面存在的偏倚，严重程度和持续时间超过相应年龄允许范围。\textbf{ADHD}：注意缺陷在行为上表现为游离于任务、缺乏持续性、难以维持注意力，多动是在不恰当的时候过多的躯体运动，起病于学龄期。\textbf{SLD}：学龄儿童在阅读、书写、拼字、表达、计算等认知学习过程中存在一种或一种以上特殊性障碍。\textbf{ASD}：以社交沟通障碍和重复刻板行为为主要特征的发育障碍性疾病。

    \item (1)智力发展正常；(2)情绪稳定而反应适度；(3)心理行为特点符合年龄；(4)人际关系的心理适应，能与人和睦相处，悦纳自己，认同他人；(5)个性的稳定和健全，表现出健康的精神风貌，客观的自我意识，行为符合道德规范，能适度耐受各种压力和应激。

    \item 童年期：思维从具体形象向抽象逻辑过渡，情感内容逐渐丰富，学校集体生活成为重要社会化途径。青春期：抽象逻辑思维占优势，情绪两极性和内隐性共存，自我意识迅猛发展，经历异性疏远期→异性爱慕期→两性初恋期，同伴影响力增强。

    \item (1)过度活动；(2)注意力不集中；(3)冲动：易兴奋、做事不顾后果、常有攻击行为；(4)学习困难、学习成绩不良；(5)有不良情绪（原发性情绪障碍）。

    \item 核心区别：ADHD有不良情绪（原发性情绪障碍），而SLD无不良情绪。ADHD主要表现为注意缺陷和多动/冲动；SLD主要表现为阅读困难、书写表达障碍或计算困难等学业技能的特殊性障碍。

    \item 核心症状：社交障碍——缺乏社交技巧，生命早期即表现出目光对视交流的缺乏。其他特征：(1)语言障碍或落后（主要就诊原因）；(2)兴趣狭隘和刻板行为；(3)认知能力障碍（大部分智力落后，部分有岛状能力）；(4)感知觉障碍：对疼痛不敏感，对某些声音或触摸表现出强烈偏好或极端反感。
\end{enumerate}
}

\newpage
